\renewcommand{\bibname}{Quellen und Anschlüsse im Projekt}
\begin{thebibliography}{99}
\addcontentsline{toc}{chapter}{Quellen und Anschlüsse im Projekt}
\small
\interlinepenalty=10000
\bibitem{projectindex}
Martin Kunze: \emph{Die Grundlagen der Mathematik}, Band~37,
Abschnitte über Index, Periode und idempotente Potenzen;
Ergänzung \emph{Index und Periode endlicher Halbgruppen}.
Projektquellen: \path{tex/b37/index-period/},
\path{tex/b37/structure/}.

\bibitem{projecttables}
Martin Kunze: \emph{Die Grundlagen der Mathematik}, Band~37,
Multiplikationstafeln und Isomorphie.
Projektquellen: \path{tex/b37/table-isomorphisms/}.

\bibitem{harju}
Tero Harju: \emph{Lecture Notes on Semigroups. Short Course}.
Universität Turku, 2010; insbesondere die Abschnitte~1.1 und~1.2.
\url{https://users.utu.fi/harju/semigroups/Semigroups2010.pdf}.

\bibitem{grillet}
Pierre A. Grillet: \emph{Commutative Semigroups}.
Kluwer, 2001, Kapitel~I, Abschnitt~5, Proposition~5.8.
Die dort beschriebene zyklische Unterhalbgruppe ist auch innerhalb
einer nichtkommutativen Halbgruppe kommutativ.
\url{https://doi.org/10.1007/978-1-4757-3389-1}.

\bibitem{pin}
Jean-Éric Pin: \emph{Mathematical Foundations of Automata Theory}.
Online-Manuskript, insbesondere Kapitel~II (Potenzen und Idempotente)
und~V (Green-Relationen und lokale Theorie).
\url{https://www.irif.fr/~jep/PDF/MPRI/MPRI.pdf}.

\bibitem{howie}
John M. Howie: \emph{Fundamentals of Semigroup Theory}.
London Mathematical Society Monographs, New Series~12,
Clarendon Press, Oxford, 1995.
\url{https://doi.org/10.1093/oso/9780198511946.001.0001}.

\bibitem{east2018}
James East, James D. Mitchell, Nik Ruškuc und Michael Torpey:
\emph{Congruence lattices of finite diagram monoids}.
Advances in Mathematics 333 (2018), 931--1003;
allgemeine Vorbereitungen, insbesondere Lemma~3.11 zum minimalen Ideal.
\url{https://doi.org/10.1016/j.aim.2018.05.016}.

\bibitem{rees1940}
David Rees: \emph{On semi-groups}.
Mathematical Proceedings of the Cambridge Philosophical Society
36 (1940), 387--400.
\url{https://doi.org/10.1017/S0305004100017436}.

\bibitem{shah2016}
M. H. Shahzamanian: \emph{The congruence eta* on semigroups}.
Quarterly Journal of Mathematics 67 (2016), 405--424;
allgemeine Vorbereitungen zu Hauptreihen und Rees-Matrixhalbgruppen.
\url{https://doi.org/10.1093/qmath/haw017}.

\bibitem{graham1968}
N. Graham, R. L. Graham und J. Rhodes:
\emph{Maximal subsemigroups of finite semigroups}.
Journal of Combinatorial Theory 4 (1968), 203--209.
\url{https://doi.org/10.1016/S0021-9800(68)80001-8}.

\bibitem{donoven2018}
C. R. Donoven, J. D. Mitchell und W. A. Wilson:
\emph{Computing maximal subsemigroups of a finite semigroup}.
Journal of Algebra 505 (2018), 559--596;
insbesondere die Lemmata~1.2 und~1.3 sowie Proposition~1.4.
\url{https://doi.org/10.1016/j.jalgebra.2018.01.044}.

\bibitem{Simon1990}
Imre Simon: \emph{Factorization forests of finite height}.
Theoretical Computer Science 72 (1990), 65--94.
\url{https://doi.org/10.1016/0304-3975(90)90047-L}.

\bibitem{branco2020}
M. J. J. Branco, G. M. S. Gomes, Jean-Éric Pin und Isabel Soler-Escrivà:
\emph{On formations of monoids}.
Journal of Pure and Applied Algebra 224 (2020), Artikel~106401;
Proposition~2.1 (für Halbgruppen).
\url{https://doi.org/10.1016/j.jpaa.2020.106401}.

\bibitem{Kufleitner2008}
Manfred Kufleitner: \emph{The Height of Factorization Forests}.
In: Mathematical Foundations of Computer Science 2008,
LNCS~5162, Springer, 2008, 443--454.
\url{https://doi.org/10.1007/978-3-540-85238-4_36}.

\bibitem{KrohnRhodes1965}
Kenneth Krohn und John Rhodes:
\emph{Algebraic theory of machines. I. Prime decomposition theorem
for finite semigroups and machines}.
Transactions of the American Mathematical Society 116 (1965), 450--464.
\url{https://doi.org/10.1090/S0002-9947-1965-0188316-1}.

\bibitem{RhodesSteinberg2009}
John Rhodes und Benjamin Steinberg:
\emph{The q-theory of Finite Semigroups}.
Springer Monographs in Mathematics, Springer, 2009;
insbesondere die Sätze~4.1.30, 4.1.45 und~4.9.1 zur Zerlegung,
zu einfachen Gruppenteilern und zur Komplexität sowie die
Abschnitte~4.12 und~4.17 zum Typ-II-Satz.
\url{https://doi.org/10.1007/978-0-387-09781-7}.

\bibitem{gillam1972}
D. W. H. Gillam, T. E. Hall und N. H. Williams:
\emph{On finite semigroups and idempotents}.
Bulletin of the London Mathematical Society 4 (1972), 143--144.
\url{https://doi.org/10.1112/blms/4.2.143}.

\bibitem{wang2016}
Guoqing Wang:
\emph{Structure of the largest idempotent-product free sequences in semigroups}.
Vorabfassung, 2016; Theorem~A und Proposition~1.1.
\url{https://arxiv.org/abs/1405.6278}.

\bibitem{hallsapir1996}
T. E. Hall und M. V. Sapir:
\emph{Idempotents, regular elements and sequences from finite semigroups}.
Discrete Mathematics 161 (1996), 151--160.
\url{https://doi.org/10.1016/0012-365X(95)00223-J}.

\bibitem{carton2007}
Olivier Carton, Dominique Perrin und Jean-Éric Pin:
\emph{Automata and semigroups recognizing infinite words}.
In: J. Flum, E. Grädel und T. Wilke (Hrsg.),
\emph{Logic and Automata: History and Perspectives},
Amsterdam University Press, 2007, 133--167; Abschnitt~3.1, Theorem~3.1.
\url{https://www.irif.fr/~jep/PDF/OmegaSem.pdf}.

\bibitem{Higgins2017}
Peter M. Higgins: \emph{Embedding in a Finite 2-Generator Semigroup}.
Glasgow Mathematical Journal 59 (2017), 61--75.
Insbesondere Satz~1.2 und S.~62 zur Einbettung mit drei idempotenten Erzeugern.
\url{https://doi.org/10.1017/S0017089516000045}.

\bibitem{Bergman2004}
George M. Bergman:
\emph{Every finite semigroup is embeddable in a finite relatively free semigroup}.
Journal of Pure and Applied Algebra 186 (2004), 1--19.
\url{https://doi.org/10.1016/S0022-4049(03)00126-9}.

\bibitem{izhakian2011}
Zur Izhakian, John Rhodes und Benjamin Steinberg:
\emph{Representation Theory of Finite Semigroups over Semirings}.
Journal of Algebra 336 (2011), 139--157; Abschnitt~2, Fact~2.1.
\url{https://doi.org/10.1016/j.jalgebra.2011.02.048}.

\bibitem{margolis2024}
Stuart Margolis, John Rhodes und Anne Schilling:
\emph{Decidability of Krohn--Rhodes complexity for all finite semigroups
and automata}. Vorabfassung, 2024; Theorem~1.5.
\url{https://arxiv.org/abs/2406.18477}.

\bibitem{margolis2025}
Stuart Margolis, John Rhodes und Anne Schilling:
\emph{Complexity of Finite Semigroups: History and Decidability}.
Semigroup Forum 111 (2025), 331--350;
insbesondere Definition~1.4 und Theorem~2.1.
\url{https://doi.org/10.1007/s00233-025-10560-z}.

\bibitem{Ash1991}
Christopher J. Ash:
\emph{Inevitable graphs: a proof of the type II conjecture and some
related decision procedures}.
International Journal of Algebra and Computation 1 (1991), 127--146.
\url{https://doi.org/10.1142/S0218196791000079}.

\bibitem{Henckell1988}
Karsten Henckell:
\emph{Pointlike sets: the finest aperiodic cover of a finite semigroup}.
Journal of Pure and Applied Algebra 55 (1988), 85--126.
\url{https://doi.org/10.1016/0022-4049(88)90042-4}.

\bibitem{VanGoolSteinberg2019}
Samuel J. van Gool und Benjamin Steinberg:
\emph{Merge Decompositions, Two-sided Krohn--Rhodes, and Aperiodic Pointlikes}.
Canadian Mathematical Bulletin 62 (2019), 199--208;
Abschnitt~4, insbesondere Theorem~4.3.
\url{https://doi.org/10.4153/CMB-2018-014-8}.

\bibitem{AlmeidaMargolisSteinbergVolkov2009}
Jorge Almeida, Stuart Margolis, Benjamin Steinberg und Mikhail Volkov:
\emph{Representation theory of finite semigroups, semigroup radicals
and formal language theory}.
Transactions of the American Mathematical Society 361 (2009), 1429--1461;
insbesondere Lemma~2.2 zur Idempotentenfaktorisierung und
Proposition~3.11 zur Hebung regulärer $\mathcal J$-Klassen.
\url{https://doi.org/10.1090/S0002-9947-08-04712-0}.

\bibitem{HenckellMargolisPinRhodes1991}
Karsten Henckell, Stuart W. Margolis, Jean-Éric Pin und John Rhodes:
\emph{Ash's type II theorem, profinite topology and Malcev products: Part I}.
International Journal of Algebra and Computation 1 (1991), 411--436;
insbesondere Satz~1.2 zu gruppenpunktartigen Teilmengen.
\url{https://doi.org/10.1142/S0218196791000298}.

\bibitem{LeSaecPinWeil1991}
Bertrand Le Saëc, Jean-Éric Pin und Pascal Weil:
\emph{Semigroups with idempotent stabilizers and applications to automata theory}.
International Journal of Algebra and Computation 1 (1991), 291--314;
Satz~1.1. Das Autorenmanuskript trägt den Titel
\emph{Finite semigroups whose stabilizers are idempotent semigroups}.
\url{https://doi.org/10.1142/S0218196791000195}.

\bibitem{DiekertWalter2016}
Volker Diekert und Tobias Walter:
\emph{Characterizing Classes of Regular Languages Using Prefix Codes of
Bounded Synchronization Delay}.
In: 43rd International Colloquium on Automata, Languages, and Programming
(ICALP 2016), LIPIcs 55, Artikel~129, 129:1--129:14;
Lemma~1 und Proposition~14 zur Rees-Zerlegung.
\url{https://doi.org/10.4230/LIPIcs.ICALP.2016.129}.

\bibitem{Henckell2004}
Karsten Henckell: \emph{Idempotent Pointlike Sets}.
International Journal of Algebra and Computation 14 (2004), 703--717.
\url{https://doi.org/10.1142/S0218196704002006}.

\bibitem{HenckellRhodesSteinberg2010StablePairs}
Karsten Henckell, John Rhodes und Benjamin Steinberg:
\emph{A Profinite Approach to Stable Pairs}.
International Journal of Algebra and Computation 20 (2010), 269--285;
Satz~4.5 zur idempotenten Punktartigkeit und Satz~6.4 zu stabilen Paaren.
\url{https://doi.org/10.1142/S0218196710005650}.

\bibitem{Rhodes1967Homomorphism}
John Rhodes: \emph{A homomorphism theorem for finite semigroups}.
Mathematical Systems Theory 1 (1967), 289--304; Abschnitt~2.
\url{https://doi.org/10.1007/BF01695164}.

\bibitem{Rhodes1970Survey}
John Rhodes: \emph{Algebraic theory of finite semigroups}.
Séminaire Dubreil. Algèbre et théorie des nombres 23 (1969--1970),
Nr.~2, Vortrag~DG10, 9~Seiten;
Proposition auf S.~DG10-07 mit den vorangehenden Definitionen.
\url{https://www.numdam.org/item/SD_1969-1970__23_2_A9_0/}.

\bibitem{DiekertKufleitnerReinhardtWalter2015}
Volker Diekert, Manfred Kufleitner, Klaus Reinhardt und Tobias Walter:
\emph{Regular Languages Are Church-Rosser Congruential}.
Journal of the ACM 62 (2015), Nr.~5, Artikel~39; Satz~5.1.
\url{https://doi.org/10.1145/2808227}.

\bibitem{DonovenKappe2023}
Casey Donoven und Luise-Charlotte Kappe:
\emph{Finite coverings of semigroups and related structures}.
International Journal of Group Theory 12 (2023), Nr.~3, 205--222;
Satz~1.4 zur Überdeckung endlicher Halbgruppen.
\url{https://doi.org/10.22108/ijgt.2022.131538.1759}.
\end{thebibliography}
